\section{Future Work}\label{sec:concl:future}

\subsection{Cross-Window Association}\label{sec:concl:future:windows}

The scores in \cref{ch:experiments} describe a window, not a recording.
Cluster indices stop at the window boundary, so the same emitter in two adjacent windows is not known to be the same object.
A cheap association is to replace each window cluster by the mean of its embeddings and to cluster those means over the recording with \gls{dbscan}.
The pulse encoder stays frozen.
A second transformer on top of those means, or of the window embeddings, could learn the association instead of leaving it to a second density step.
Either construction yields one partition per scenario.

\subsection{Streaming Inference}\label{sec:concl:future:online}

The encoder is bidirectional over a block of $L$ pulses.
A live stream does not arrive as a block.
Key--value (KV) caching stores the keys and values of pulses already in the window, so a new pulse attends against that cache instead of forcing a full forward pass.
A mask then chooses which cached pulses remain visible: a sliding window of length $L$, or a causal band if later pulses must not leak into the present.
That is incremental inference.
Updating the weights from the stream is a separate problem and is not required for the cache.

\subsection{Capacity}\label{sec:concl:future:capacity}

The encoder may be oversized for this corpus.
A shared trunk of four layers and two four-layer branches, each of width $256$ with inner feed-forward dimension $2048$, give about sixteen million parameters (\cref{tab:exp:efficiency,sec:method:arch}), while most windows already sit at the score ceiling.
A narrower or shallower stack, scored on the same split, would test whether those scores survive a smaller model.
Width, depth, and feed-forward expansion can be reduced independently.
The useful product is the smallest encoder that still matches the tables of \cref{ch:experiments}.

\subsection{Hyperparameters}\label{sec:concl:future:hparams}

Joint training uses the unweighted sum $\Loss_{\mathrm{joint}} = \Loss_{\mathrm{em}} + \Loss_{\mathrm{md}}$ (\cref{eq:method:loss_total}).
A relative weight between the two terms would show whether the emitter cost of the joint loss can be reduced without giving the mode branch back to merging or splitting.
The ratio need not be fixed: weights that follow the running values of $\Loss_{\mathrm{em}}$ and $\Loss_{\mathrm{md}}$ could keep the two terms on a comparable scale as training proceeds.

The rotary scales were a single unlearned choice, $\gamma_t=\gamma_x=\gamma_y=10$, far below the window length $L$, and the rotary base was left at $b=10^{4}$ (\cref{sec:method:attention:rope}).
RoPE-TOA already sits at the emitter ceiling of \cref{tab:exp:attn_em}, so retuning $\gamma_t$ is unlikely to move that score.
RoPE-az still has headroom, and tying $\gamma_x$ and $\gamma_y$ to $\gamma_t$ is one candidate reason it does not match TOA-only rotation.
Window-local \gls{toa} and planar incidence both occupy $[0,1]$ after \cref{sec:method:data:norm}, yet a numerical gap of a given size is not the same physical distance in the two coordinates, and the deinterleaving split gives each incidence axis only half the rotation planes of a TOA head.
A larger $\gamma_x$ and $\gamma_y$ would give nearby platforms more phase; a smaller pair would keep phases unique across a wide scene and would flatten those nearby distinctions.
A per-coordinate sweep, or treating $\gamma$ as a learned scalar like $\lambda$, would show whether the shared setting, rather than incidence already being an input feature, is what left RoPE-az short of RoPE-TOA.

\subsection{Mode Classification Head}\label{sec:concl:future:cls}

Mode labels are already global.
After contrastive training, a classification head can be attached to the mode-branch tokens $\vect{h}^{\mathrm{md}}_n$ of \cref{sec:method:arch}, before the projection that produces $\vect{z}^{\mathrm{md}}_n$.
The encoder stays frozen; only the head is trained to name types from the training catalogue.
Mode clustering already recovers those types on almost every window (\cref{tab:exp:attn_md}), so the branch representation already separates them on the manifold.
A linear head on the same tokens tests whether that separation is readable as class logits.
Clustering remains the decoder for types that are not in the catalogue.
The missing single-task counterpart is a classical mode clusterer or classifier on the trains of an emitter-only encoder.

\subsection{Dataset Hardness}\label{sec:concl:future:hardness}

Mean $\ARI$ hides which recordings are hard.
Preliminary checks on this corpus used label entropy, the consecutive-label transition rate, and a silhouette score on the raw \gls{pdw} features; none of them tracked per-window $\ARI$.
Entropy and related information-theoretic summaries treat pulses as exchangeable draws, so they miss both class imbalance and the temporal and spatial dependence that make a mixture hard.
Those checks also omitted the mutual information between the emitter and mode labels of a window (\cref{eq:bg:mi_def}).
That quantity would say how much one ground-truth partition determines the other, which was not measured here.
A more useful scale is the error of a simple baseline that has to solve a related task on the same window: a histogram classifier on carrier frequency, pulse width, and \gls{pri} for modes, or Feature \gls{dbscan} for either partition.
Windows where that baseline fails are the hard ones, and they give a comparison axis for a later generator or for operational data.

\subsection{Held-Out Types and Operational Data}\label{sec:concl:future:shift}

Held-out modulation types are the experiment that the clustering decoder is for.
Until that split is scored, the mode branch is a grouping of seen catalogue types.
Operational recordings, and simulator drop or spurious-pulse regimes as an actual test set, are the corresponding check on distribution shift.
